\documentclass[12pt,a4paper]{article}

\usepackage[utf8]{inputenc}
\usepackage[T1]{fontenc}
\usepackage{amsmath, amssymb, amsthm}
\usepackage[dvipsnames]{xcolor}
\usepackage{geometry}
\usepackage{booktabs}
\usepackage{hyperref}
\usepackage{parskip}

\geometry{margin=2.5cm}

\hypersetup{colorlinks=true, linkcolor=blue!60!black,
            citecolor=blue!60!black, urlcolor=blue!60!black}

\newtheorem{remark}{Observacion}

\title{\textbf{Cuadratura mediante Reciprocidad Dual en 1D}\\[0.5em]
\large Analogia con el Caso 2D y Fundamentos del Metodo\\[0.3em]
\normalsize Notas para un Curso de Maestria en Boundary Element Method}
\author{W.F. Florez\\
\small Wessex Institute of Technology}
\date{}

\begin{document}

\maketitle
\tableofcontents
\newpage

% ─────────────────────────────────────────────────────────────
\section{Motivacion: El Caso 1D como Punto de Partida}
% ─────────────────────────────────────────────────────────────

Antes de abordar la conversion de integrales de dominio en 2D mediante el
metodo de soluciones particulares, es instructivo estudiar el caso
unidimensional. En 1D toda la maquinaria del metodo se reduce a sus elementos
esenciales, sin la complejidad geometrica de la frontera ni la necesidad de
cuadratura sobre elementos. Esto permite aislar y estudiar de forma pura el
error de colocacion RBF, que es una de las dos fuentes de error presentes
en el caso 2D.

% ─────────────────────────────────────────────────────────────
\section{Planteamiento del Problema}
% ─────────────────────────────────────────────────────────────

Se desea evaluar la integral:

\begin{equation}
    I = \int_a^b g(x)\, dx
    \label{eq:integral_1d}
\end{equation}

La estrategia es aproximar $g(x)$ mediante una combinacion lineal de
funciones base $\{f_i(x)\}$ evaluadas en $N$ nodos de colocacion
$\{x_i\}_{i=1}^N \subset (a,b)$:

\begin{equation}
    g(x) \approx \sum_{i=1}^{N} \alpha_i f_i(x)
    \label{eq:aprox_g}
\end{equation}

y convertir la integral de dominio en una evaluacion en los extremos
$\{a, b\}$ (la ``frontera'' del intervalo en 1D).

% ─────────────────────────────────────────────────────────────
\section{Conversion de la Integral: el Teorema Fundamental del Calculo}
% ─────────────────────────────────────────────────────────────

\subsection{Condicion sobre las Funciones Base}

Se exige que cada $f_i(x)$ sea la derivada de alguna funcion $\phi_i(x)$:

\begin{equation}
    \frac{d\phi_i(x)}{dx} = f_i(x)
    \label{eq:condicion_phi}
\end{equation}

Esta es la condicion analoga a $\nabla^2\hat{\phi}_i = f_i(r)$ en el caso 2D:
se fuerza que la funcion base sea la imagen de un operador diferencial que
admite una identidad integral.

\subsection{Conversion}

Sustituyendo \eqref{eq:aprox_g} y \eqref{eq:condicion_phi} en
\eqref{eq:integral_1d}:

\begin{align}
    I &= \int_a^b g(x)\,dx
     \approx \sum_{i=1}^{N} \alpha_i \int_a^b f_i(x)\,dx
     = \sum_{i=1}^{N} \alpha_i \int_a^b \frac{d\phi_i}{dx}\,dx
\end{align}

Aplicando el \textbf{Teorema Fundamental del Calculo}:

\begin{equation}
    \int_a^b \frac{d\phi_i}{dx}\,dx = \left[\phi_i(x)\right]_a^b
    = \phi_i(b) - \phi_i(a)
\end{equation}

y el resultado final es:

\begin{equation}
\boxed{
    I \approx \sum_{i=1}^{N} \alpha_i \bigl(\phi_i(b) - \phi_i(a)\bigr)
}
\label{eq:resultado_1d}
\end{equation}

La integral de dominio queda expresada como una \textbf{evaluacion puntual
en los extremos del intervalo}, sin ninguna cuadratura adicional.

\begin{remark}
El Teorema Fundamental del Calculo es precisamente el Teorema de la
Divergencia en 1D. La ``frontera'' del intervalo $(a,b)$ son los dos
puntos $\{a, b\}$, y la ``integral de frontera''
$\int_\Gamma \frac{\partial\hat\phi}{\partial n}\,d\Gamma$ se reduce a
$\phi(b)\cdot(+1) + \phi(a)\cdot(-1) = \phi(b) - \phi(a)$,
donde $\pm 1$ son las normales exteriores en cada extremo.
\end{remark}

% ─────────────────────────────────────────────────────────────
\section{Tabla Comparativa: 1D vs 2D}
% ─────────────────────────────────────────────────────────────

\begin{center}
\begin{tabular}{lll}
\toprule
Elemento & 1D & 2D \\
\midrule
Integral de dominio
    & $\displaystyle\int_a^b g\,dx$
    & $\displaystyle\int_\Omega b\,d\Omega$ \\[1em]
Condicion sobre base
    & $\dfrac{d\phi_i}{dx} = f_i$
    & $\nabla^2\hat\phi_i = f_i(r)$ \\[1em]
Identidad integral
    & Teorema Fundamental
    & Teorema de la Divergencia \\[0.5em]
``Frontera''
    & Dos puntos $\{a,\,b\}$
    & Curva $\Gamma$ \\[0.5em]
Evaluacion en frontera
    & $\phi_i(b) - \phi_i(a)$
    & $\displaystyle\int_\Gamma \frac{\partial\hat\phi_i}{\partial n}\,d\Gamma$ \\[1em]
Cuadratura sobre $\Gamma$
    & No necesaria
    & Cuadratura de Gauss por elemento \\[0.5em]
Fuentes de error
    & Solo colocacion RBF
    & Geometria + colocacion RBF \\
\bottomrule
\end{tabular}
\end{center}

% ─────────────────────────────────────────────────────────────
\section{Obtencion de $\phi_i(x)$ para la RBF $f_i = 1 + |x - x_i|^3$}
% ─────────────────────────────────────────────────────────────

\subsection{Integracion Directa}

Se elige la funcion base radial 1D:

\begin{equation}
    f_i(x) = 1 + |x - x_i|^3, \qquad r_i = x - x_i
\end{equation}

Hay que resolver $\dfrac{d\phi_i}{dx} = 1 + |r_i|^3$ integrando directamente.
El termino constante integra trivialmente. Para el termino $|r_i|^3$,
observamos que:

\begin{equation}
    \frac{d}{dx}\bigl[r_i\,|r_i|^3\bigr]
    = |r_i|^3 + r_i \cdot 3|r_i|^2\,\text{sgn}(r_i)
    = |r_i|^3 + 3|r_i|^3 = 4|r_i|^3
\end{equation}

por tanto:

\begin{equation}
    \int |r_i|^3\,dx = \frac{r_i\,|r_i|^3}{4} + C
    = \frac{(x-x_i)|x-x_i|^3}{4} + C
\end{equation}

\subsection{Resultado}

Tomando la constante de integracion $C = 0$:

\begin{equation}
\boxed{
    \phi_i(x) = x + \frac{(x - x_i)\,|x - x_i|^3}{4}
}
\label{eq:phi_1d}
\end{equation}

\subsection{Verificacion}

Derivando \eqref{eq:phi_1d}:

\begin{align}
    \frac{d\phi_i}{dx}
    &= 1 + \frac{d}{dx}\left[\frac{(x-x_i)|x-x_i|^3}{4}\right] \notag \\
    &= 1 + \frac{|x-x_i|^3 + (x-x_i)\cdot 3|x-x_i|^2\,\text{sgn}(x-x_i)}{4} \notag \\
    &= 1 + \frac{|x-x_i|^3 + 3|x-x_i|^3}{4} \notag \\
    &= 1 + |x-x_i|^3 = f_i(x) \qquad \checkmark
\end{align}

% ─────────────────────────────────────────────────────────────
\section{Determinacion de los Coeficientes $\alpha_i$ por Colocacion}
% ─────────────────────────────────────────────────────────────

Se conoce $g(x)$ en los $N$ nodos interiores $\{x_i\}_{i=1}^N \subset (a,b)$.
Evaluando la aproximacion \eqref{eq:aprox_g} en cada nodo:

\begin{equation}
    g(x_k) = \sum_{j=1}^{N} \alpha_j f_j(x_k)
    = \sum_{j=1}^{N} \alpha_j \bigl(1 + |x_k - x_j|^3\bigr),
    \quad k = 1,\ldots,N
\end{equation}

lo que genera el sistema lineal:

\begin{equation}
    \mathbf{F}\,\boldsymbol{\alpha} = \mathbf{g},
    \qquad F_{kj} = 1 + |x_k - x_j|^3,
    \qquad g_k = g(x_k)
\end{equation}

con solucion:

\begin{equation}
    \boldsymbol{\alpha} = \mathbf{F}^{-1}\mathbf{g}
\end{equation}

\begin{remark}
Los nodos de colocacion deben ser \textbf{interiores}: $x_i \in (a,b)$,
nunca en los extremos $a$ o $b$. Si $x_i = a$ o $x_i = b$, entonces
$\phi_i(a)$ o $\phi_i(b)$ involucran $r_i = 0$, lo que no causa singularidad
en 1D (a diferencia del caso 2D), pero si situa un nodo fuente sobre la
``frontera'', lo cual es equivalente al problema de near-singularity en 2D.
\end{remark}

% ─────────────────────────────────────────────────────────────
\section{Algoritmo Completo}
% ─────────────────────────────────────────────────────────────

El procedimiento completo para evaluar $I = \int_a^b g(x)\,dx$ es:

\begin{enumerate}
    \item Elegir $N$ nodos de colocacion interiores
    $\{x_i\}_{i=1}^N \subset (a,b)$.

    \item Evaluar $g(x_i)$ en cada nodo para construir $\mathbf{g}$.

    \item Construir la matriz de colocacion $F_{ij} = 1 + |x_i - x_j|^3$.

    \item Resolver $\mathbf{F}\boldsymbol{\alpha} = \mathbf{g}$.

    \item Evaluar $\phi_i(b) - \phi_i(a)$ para cada $i$ usando
    \eqref{eq:phi_1d}.

    \item Calcular $I \approx \sum_{i=1}^N \alpha_i\bigl(\phi_i(b)
    - \phi_i(a)\bigr)$.
\end{enumerate}

% ─────────────────────────────────────────────────────────────
\section{Resultados Numericos}
% ─────────────────────────────────────────────────────────────

\subsection{Ejemplo Principal: $g(x) = x^2$ en $[0,1]$}

La solucion analitica es:

\begin{equation}
    I = \int_0^1 x^2\,dx = \frac{1}{3}
\end{equation}

\subsubsection{Estudio de Convergencia}

Se usan $N$ nodos equidistantes interiores en $(0,1)$ y se registra el
error relativo y el condicionamiento de $\mathbf{F}$:

\begin{center}
\begin{tabular}{rrrr}
\toprule
$N$ & $I_{\text{DR}}$ & Error (\%) & $\mathrm{cond}(\mathbf{F})$ \\
\midrule
 2 & 0.2870370370 & 13.889  &         55 \\
 3 & 0.3196358268 &  4.109  &        155 \\
 5 & 0.3290538035 &  1.284  &      2\,715 \\
 8 & 0.3320184170 &  0.394  &     20\,518 \\
12 & 0.3328809512 &  0.136  &    105\,141 \\
20 & 0.3332212606 &  0.034  &    785\,205 \\
\bottomrule
\multicolumn{4}{r}{$I_{\text{exacta}} = 0.3333333333$}
\end{tabular}
\end{center}

El error decrece monotonamente con $N$, pero el condicionamiento crece
rapidamente --- de 55 con $N=2$ a mas de $10^5$ con $N=12$.

\subsubsection{Tension entre Precision y Condicionamiento}

Este resultado pone de manifiesto una tension fundamental en el metodo:
aumentar $N$ mejora la aproximacion de $g$ pero deteriora el
condicionamiento de $\mathbf{F}$. En 1D esta tension es visible de forma
pura porque no hay error geometrico ni de cuadratura que la enmascare.
En la practica existe un $N$ optimo que minimiza el error total teniendo
en cuenta tanto la aproximacion RBF como la propagacion de errores de
redondeo a traves de $\mathbf{F}^{-1}$.

\subsection{Prueba con Distintas Funciones $g(x)$, $N=8$}

\begin{center}
\begin{tabular}{lrrl}
\toprule
$g(x)$ & $I_{\text{exacta}}$ & $I_{\text{DR}}$ & Error (\%) \\
\midrule
$x^2$       & 0.33333333 & 0.33201842 & 0.394 \\
$x^3$       & 0.25000000 & 0.24750286 & 0.999 \\
$x^4$       & 0.20000000 & 0.19572460 & 2.138 \\
$\sin(\pi x)$ & 0.63661977 & 0.64193920 & 0.836 \\
$e^x$       & 1.71828183 & 1.72011136 & 0.106 \\
\bottomrule
\end{tabular}
\end{center}

Se observa que el error crece con el grado del polinomio: la RBF
$f = 1 + |r|^3$ aproxima mejor funciones de bajo grado, consistente con
el hecho de que el spline cubico $|r|^3$ reproduce exactamente polinomios
hasta grado 3 pero no mas alla.

% ─────────────────────────────────────────────────────────────
\section{Observaciones Pedagogicas}
% ─────────────────────────────────────────────────────────────

\subsection{Aislamiento del Error de Colocacion RBF}

En 1D el unico error presente es el de colocacion RBF. No hay:
\begin{itemize}
    \item error geometrico (la ``frontera'' $\{a,b\}$ es exacta),
    \item error de cuadratura sobre $\Gamma$ (la evaluacion en $a$ y $b$
    es exacta).
\end{itemize}

Esto hace del caso 1D el laboratorio ideal para estudiar el error de
aproximacion RBF de forma pura, antes de introducir las fuentes adicionales
de error presentes en 2D.

\subsection{Secuencia Pedagogica Recomendada}

La progresion natural para el curso es:

\begin{enumerate}
    \item \textbf{Caso 1D:} motivacion limpia, sin geometria.
    Error unico: colocacion RBF. La ``cuadratura de frontera'' es
    simplemente $\phi_i(b) - \phi_i(a)$.

    \item \textbf{Caso 2D, dominio cuadrado:} geometria exacta con
    elementos constantes. Aparece la cuadratura de Gauss sobre $\Gamma$,
    pero sin error geometrico. Las dos fuentes de error son: RBF y Gauss.

    \item \textbf{Caso 2D, dominio circular:} geometria aproximada por
    cuerdas. Las tres fuentes de error estan presentes: RBF, Gauss, y
    geometria. Se estudia la necesidad de refinar ambas discretizaciones
    en forma coordinada.
\end{enumerate}

\subsection{Preguntas Propuestas para el Estudiante}

\begin{enumerate}
    \item Mostrar que la eleccion de la constante de integracion $C$
    en \eqref{eq:phi_1d} no afecta el resultado \eqref{eq:resultado_1d}.

    \item Derivar $\phi_i(x)$ para $f_i(x) = e^{-\beta(x-x_i)^2}$
    (RBF Gaussiana 1D) y discutir si la integral puede obtenerse en
    forma cerrada.

    \item Para $g(x) = x^n$ con $n$ entero, determinar cuantos nodos
    de colocacion $N$ son necesarios para que el error sea menor que
    $10^{-4}$, y graficar la tension error-condicionamiento.

    \item Comparar el resultado de la cuadratura DR con la regla de
    Simpson y la cuadratura de Gauss para el mismo numero de evaluaciones
    de $g(x)$. ?`Cuando es ventajosa la estrategia DR?
\end{enumerate}

\end{document}